\paragraph{稳定负子空间、端点谱迹与时间纤维：Toeplitz 否证的可编译接口}
本段将“有限否证器”的负谱结构进一步压缩为可复算的模型空间截断像与端点读出不变量，并把 window-\(6\) 的最小热化闭合、Witt--Fourier 双桥与同步核中心六重聚合纳入同一套自动编号的结论口径。

\begin{conclusion}[Toeplitz 否证器的可编译性：稳定负子空间是模型空间的截断像]\label{con:xi-toeplitz-compiled-model-space-truncation}
在有限 Blaschke 缺陷设定下，令 \(\kappa:=\deg(B_{F_{\zeta}})=\dim K_{B_{F_{\zeta}}}\)。
则存在 \(N_0\) 与一族随 \(N\) 稳定的显式线性嵌入
\[
\iota_N:K_{B_{F_{\zeta}}}\hookrightarrow \CC^{N+1}\qquad(N\ge N_0)
\]
（例如取 Hardy 基 \(\{1,w,w^2,\dots\}\) 的前 \(N{+}1\) 个 Taylor 系数坐标），使 Toeplitz 二次型
\(
x\mapsto x^\ast\mathbf{T}_N x
\)
在 \(\iota_N(K_{B_{F_{\zeta}}})\) 上的限制具有\textbf{恰好 \(\kappa\) 维负定部分}，且该负定部分在 \(N\to\infty\) 的稳定区间内保持不变。
因此一旦观测到 \(\mathbf{T}_N\) 的稳定负子空间，即可在有限维上实现一个与压缩移位 \(A_{B_{F_\zeta}}\) 的点谱残差轴同型的“残差轴”读出，并以其点谱数据反演圆盘内缺陷点（按重数计）。
\end{conclusion}
\begin{proof}[证明要点]
稳定嵌入与“负子空间 \(=\) 模型空间截断像”的同型性见定理 \ref{thm:xi-toeplitz-witness-generated-by-model-space}；
点谱残差轴与 Blaschke 零点的对应见定理 \ref{thm:blaschke-point-spectrum-correspondence}。
结合谱读出接口（例如定理 \ref{thm:xi-resonance-pencil-reconstruction} 及结论 \ref{con:xi-finite-falsifier-offline-fiber-readout}）即可把稳定负子空间提升为有限维可审计的缺陷位置反演。
\end{proof}

\begin{conclusion}[失败见证的极小生成律：否证由点谱残差轴强制生成]\label{con:xi-toeplitz-witness-minimal-generator-axis}
沿用上条结论的有限缺陷设定。
当 \(N\) 充分大时，Toeplitz 层级的一切负方向都可由点谱残差轴
\(
\mathcal H_{\mathrm{pp}}(A_{B_{F_\zeta}})=K_{B_{F_\zeta}}
\)
的点模态经 \(\iota_N\) 的像生成；从而 Toeplitz 否证的“极小生成维数”严格等于
\(
\deg(B_{F_\zeta})=\kappa
\)。
因此有限否证器是可审计的有限维生成对象，而非数值不稳定的弥散副产物。
\end{conclusion}
\begin{proof}[证明要点]
负惯性指标在高阶截断上的稳定性见定理 \ref{thm:xi-toeplitz-inertia-stabilizes-to-kappa}；
点模态生成性与极小维数结论由定理 \ref{thm:xi-toeplitz-witness-generated-by-model-space} 直接给出。
\end{proof}

\begin{conclusion}[端点通量的算子化同一：Poisson--Busemann 权重等于 Julia 指标与端点评价算子的迹]\label{con:xi-endpoint-flux-operator-trace-identity}
设 \(B\) 为有限 Blaschke 乘积，零点多重集为 \(\{a_k\}_{k=1}^{\kappa}\subset\mathbb{D}\)，并记端点 Poisson 权重
\[
P_a(-1)=\frac{1-|a|^2}{|1+a|^2}.
\]
定义端点通量（吸收系数）
\[
\Phi_{-1}(B):=\sum_{k=1}^{\kappa}P_{a_k}(-1),
\]
并令 \(K_B:=H^2\ominus BH^2\)、端点评价泛函 \(L_{-1}(f):=f(-1)\)。
则有严格恒等式
\[
\boxed{\ \Phi_{-1}(B)=\mathcal{J}_{-1}(B)=\mathrm{tr}\!\bigl(L_{-1}^\ast L_{-1}\bigr)\ },
\]
从而端点吸收强度被提升为“端点评价算子”的可谱化迹不变量，并与点模态边界耦合平方和完全等价。
\end{conclusion}
\begin{proof}[证明要点]
这是定理 \ref{thm:xi-endpoint-flux-pointmode-coupling} 的结论（并用到 \eqref{eq:app-poisson-kernel-endpoint-minus1} 的 Poisson 端点闭式与 Julia 指标接口）。
\end{proof}

\begin{conclusion}[Artin 分解下的端点通量可加性：通道归一化行列式因子的 Poisson 计数律]\label{con:xi-endpoint-flux-artin-additivity}
设 \(G\) 为有限群，并沿用命题 \ref{prop:ihara-artin-factorization} 的群扩张 Hashimoto 记号。
固定 \(u>0\)，记 \(\lambda_1(u):=\rho(B_{\mathbf 1}(u))\)，并对每个不可约表示 \(\varrho\in\widehat G\) 定义归一化行列式因子
\[
F_\varrho(t):=\det\!\Bigl(I-(t/\lambda_1(u))B_\varrho(u)\Bigr).
\]
若对每个 \(\varrho\)，\(F_\varrho\) 在单位圆上为 inner（等价地：其零点全部位于 \(\mathbb D\) 且边界模长恒为 \(1\)，从而 \(F_\varrho\)（除去单位模常数后）为有限 Blaschke 乘积），则端点通量满足严格表示论可加性：
\[
\boxed{\;
\Phi_{-1}\!\Bigl(\prod_{\varrho\in\widehat G} F_\varrho(t)^{\dim\varrho}\Bigr)
=
\sum_{\varrho\in\widehat G}(\dim\varrho)\,\Phi_{-1}(F_\varrho)\; }.
\]
在 \(G=S_4\) 且取 \(\varrho=\mathrm{sgn}\) 的情形，记
\[
R_{\det}(u,t):=\frac{F_{\mathrm{sgn}}(t)}{F_{\mathbf 1}(t)}.
\]
则可定义跨支端点通量差
\[
\Phi_{-1}(R_{\det}):=\Phi_{-1}(F_{\mathrm{sgn}})-\Phi_{-1}(F_{\mathbf 1}),
\]
其由零点 Poisson 权重完全确定，并可借由结论 \ref{con:xi-endpoint-flux-operator-trace-identity} 的迹表示提升为可谱化不变量。
\end{conclusion}
\begin{proof}[证明要点]
inner 假设把每个 \(F_\varrho\)（除去单位模常数后）识别为有限 Blaschke 乘积，其零点多重集按乘积并合且幂次引入重数因子。
由定义 \(\Phi_{-1}=\sum_k P_{a_k}(-1)\) 对乘积严格可加，立即得到所示恒等式；\(\Phi_{-1}(R_{\det})\) 为两通道通量之差。证毕。
\end{proof}

\begin{conclusion}[端点 Poisson 通量的可审计迹公式：\texorpdfstring{$t=-1$}{t=-1} 处 resolvent 的通道差]\label{con:xi-endpoint-flux-resolvent-trace-channel-difference}
令 \(u\mapsto M^{\pm}(u)\) 为某个复邻域内解析的有限维算子族，并假设 \(I+M^{\pm}(u)\) 可逆。
定义行列式比值
\[
R(u,t):=\frac{\det(I-tM^{+}(u))}{\det(I-tM^{-}(u))}.
\]
若进一步 \(R(u,\cdot)\) 在 \(t=-1\) 处解析且 \(R(u,-1)\neq 0\)，则可定义端点通量
\[
\Phi_{-1}(R(u,\cdot))
:=
-\Re\left(\frac{\dd}{\dd t}\log R(u,t)\Big|_{t=-1}\right).
\]
此时成立严格迹恒等式
\[
\boxed{\;
\Phi_{-1}(R(u,\cdot))
=
\Re\,\mathrm{tr}\!\Bigl((I+M^{+}(u))^{-1}M^{+}(u)-(I+M^{-}(u))^{-1}M^{-}(u)\Bigr)\; }.
\]
进一步，若存在整数 \(A_{\mathrm{flip}}\ge 1\) 以及算子 \(M_0,B\) 使得在 \(u\downarrow 0\) 下有
\[
M^{\pm}(u)=M_0\pm u^{A_{\mathrm{flip}}}B+O\bigl(u^{A_{\mathrm{flip}}+1}\bigr),
\qquad
\det(I+M_0)\neq 0,
\]
则端点通量差具有首项展开
\[
\boxed{\;
\Phi_{-1}(R(u,\cdot))
=
2u^{A_{\mathrm{flip}}}\Re\,\mathrm{tr}\!\Bigl((I+M_0)^{-1}B(I+M_0)^{-1}\Bigr)
+O\bigl(u^{A_{\mathrm{flip}}+1}\bigr)\; }.
\]
特别地，若 \(\Re\,\mathrm{tr}\!\bigl((I+M_0)^{-1}B(I+M_0)^{-1}\bigr)\neq 0\)，则 \(\operatorname{ord}_u \Phi_{-1}(R(u,\cdot))=A_{\mathrm{flip}}\)。
\end{conclusion}
\begin{proof}[证明要点]
由恒等式
\(
\frac{\dd}{\dd t}\log\det(I-tM)=-\mathrm{tr}\!\bigl((I-tM)^{-1}M\bigr)
\)
得
\[
\frac{\dd}{\dd t}\log R(u,t)
=
-\mathrm{tr}\!\bigl((I-tM^{+}(u))^{-1}M^{+}(u)\bigr)
+\mathrm{tr}\!\bigl((I-tM^{-}(u))^{-1}M^{-}(u)\bigr),
\]
代入 \(t=-1\) 并取实部即得第一式。
\par
对首项展开，令 \(K:=(I+M_0)^{-1}\)。
由 Fr\'echet 展开
\(
(I+M_0\pm u^{A_{\mathrm{flip}}}B)^{-1}=K\mp u^{A_{\mathrm{flip}}}KBK+O(u^{A_{\mathrm{flip}}+1})
\)，
并用恒等式 \((I+M)^{-1}M=I-(I+M)^{-1}\)，
得到
\[
(I+M^{+}(u))^{-1}M^{+}(u)-(I+M^{-}(u))^{-1}M^{-}(u)
=(I+M^{-}(u))^{-1}-(I+M^{+}(u))^{-1}
=2u^{A_{\mathrm{flip}}}KBK+O(u^{A_{\mathrm{flip}}+1}),
\]
取迹并取实部即得所示首项系数。
\end{proof}

\begin{conclusion}[端点通量的两条稀疏证书：深度计数上界与端点读出能量预算]\label{con:xi-endpoint-flux-sparsity-and-readout-budget}
沿用结论 \ref{con:xi-endpoint-flux-operator-trace-identity} 的有限 Blaschke 设定。
对任意阈值 \(\delta>0\) 定义端点深度计数
\[
N_\delta(B):=\#\{1\le k\le \kappa:\ P_{a_k}(-1)\ge \delta\}.
\]
则有严格上界
\[
\boxed{\ N_\delta(B)\ \le\ \frac{\Phi_{-1}(B)}{\delta}\ }.
\]
此外，对任意 \(f\in K_B\) 有端点读出能量预算
\[
\boxed{\ |f(-1)|^2\ \le\ \Phi_{-1}(B)\,\|f\|^2\ }.
\]
并且常数 \(\Phi_{-1}(B)\) 为最优：存在 \(f_\star\in K_B\setminus\{0\}\) 使等号成立。
\end{conclusion}
\begin{proof}[证明要点]
第一条由定义直接得到：若 \(P_{a_k}(-1)\ge \delta\) 的项有 \(N_\delta(B)\) 个，则
\(
\Phi_{-1}(B)=\sum_k P_{a_k}(-1)\ge N_\delta(B)\delta
\)，从而 \(N_\delta(B)\le \Phi_{-1}(B)/\delta\)。
\par
第二条由 Cauchy--Schwarz 给出：\(|f(-1)|\le \|L_{-1}\|\,\|f\|\)。
在有限 Blaschke 设定下 \(\dim K_B=\kappa<\infty\)，算子 \(L_{-1}^\ast L_{-1}\) 为正的秩一算子，故
\[
\|L_{-1}\|^2=\mathrm{tr}(L_{-1}^\ast L_{-1})=\Phi_{-1}(B)
\]
（结论 \ref{con:xi-endpoint-flux-operator-trace-identity}）。
代回即可得 \(|f(-1)|^2\le \Phi_{-1}(B)\|f\|^2\)。
最优性由再生核向量的取等条件（或等价的秩一算子谱分解）给出。
\end{proof}

\begin{conclusion}[时间纤维的内生性：上半平面虚部即端点 \texorpdfstring{$-1$}{-1} 的 Poisson--Busemann 密度]\label{con:xi-timefiber-poisson-busemann-identity}
令 \(t=x+\mathrm{i}y\in\mathbb{H}\)，并取 Cayley 变换
\[
w=\mathcal{C}(t)=\frac{1+\mathrm{i}t}{1-\mathrm{i}t}\in\mathbb{D}.
\]
则成立严格恒等式
\[
\boxed{\ y=B^{-1}(w)=\frac{1-|w|^2}{|1+w|^2}\ }.
\]
因此在端点紧化口径下，“时间纤维坐标”并非外加连续轴，而被严格同一为端点 \(-1\) 的谐测度密度（Poisson--Busemann 权重）；高高度信息被强制压缩为 \(w\approx-1\) 的边界层贡献，从而形成可核验的端点局部化原则。
\end{conclusion}
\begin{proof}[证明要点]
见推论 \ref{cor:dephys-poisson-busemann-timefiber}（并引用命题 \ref{prop:app-busemann-poisson-minus1} 的 Busemann/Poisson 定义接口）。
\end{proof}

\paragraph{window-\texorpdfstring{$6$}{6} 的最小热化闭合：推送核、谱隙与 lumpability 障碍}

\begin{conclusion}[window-\texorpdfstring{$6$}{6} 的最小热化闭合：推送核的可逆性、细致平衡与显式谱隙]\label{con:xi-window6-p6-minimal-thermalization-closure}
以推送邻接计数 \(A_6\) 定义 \(X_6\) 上的推送核 \(P_6\) 与平稳分布 \(\pi_6\)（生成片段 \eqref{eq:window6_pushforward_markov_kernel}）：
\[
\pi_6(x)=\frac{d^{\mathrm{bin}}_6(x)}{64},\qquad
P_6(x,x')=\frac{A_6(x,x')}{6\,d^{\mathrm{bin}}_6(x)}.
\]
则 \(P_6\) 满足细致平衡
\(
\pi_6(x)P_6(x,x')=\pi_6(x')P_6(x',x)
\)，
因而可逆并以 \(\pi_6\) 为平稳分布。
谱隙审计给出可复核常数（生成片段 \eqref{eq:window6_pushforward_markov_spectral_gap}）
\[
\lambda_2(P_6)\approx 0.4841207858,\qquad
\lambda_\star(P_6)\approx 0.6030939754,
\]
从而对任意初始分布 \(\mu\) 存在常数 \(C(\mu)\) 使对所有 \(t\ge 0\) 有指数混合上界
\[
\|\mu P_6^t-\pi_6\|_{\mathrm{TV}}\le C(\mu)\,\lambda_\star(P_6)^t.
\]
\end{conclusion}
\begin{proof}[证明要点]
核与细致平衡见定理 \ref{thm:terminal-foldbin6-pushforward-markov} 与 \eqref{eq:window6_pushforward_markov_kernel}；
谱隙常数与 TV 上界见推论 \ref{cor:terminal-foldbin6-minimal-markov-closure-gap} 及 \eqref{eq:window6_pushforward_markov_spectral_gap}。
\end{proof}

\begin{conclusion}[从谱隙到混合时间的硬常数：window-\texorpdfstring{$6$}{6} 的热化时间尺度为 \texorpdfstring{$1/(-\log\lambda_\star)$}{1/(-log lambda*)}]\label{con:xi-window6-mixing-time-hard-constant}
沿用上条结论的记号。
若对给定初始分布 \(\mu\) 定义收敛时间
\[
t(\varepsilon;\mu):=\min\bigl\{t\ge 0:\ \|\mu P_6^t-\pi_6\|_{\mathrm{TV}}\le \varepsilon\bigr\},
\]
则由 \(\|\mu P_6^t-\pi_6\|_{\mathrm{TV}}\le C(\mu)\lambda_\star(P_6)^t\) 得到显式上界
\[
\boxed{\ t(\varepsilon;\mu)\ \le\ \frac{\log(C(\mu)/\varepsilon)}{-\log\lambda_\star(P_6)}\ }
\approx 1.9775264\cdot \log\!\bigl(C(\mu)/\varepsilon\bigr).
\]
因此 window-\(6\) 的无记忆热化闭合在时间尺度上携带一个实现无关的协议级硬常数 \(1/(-\log\lambda_\star(P_6))\)，并由谱证书直接输出而不依赖额外拟合。
\end{conclusion}
\begin{proof}[证明要点]
由推论 \ref{cor:terminal-foldbin6-minimal-markov-closure-gap} 的 TV 指数收敛式，取对数即得；
数值常数由 \eqref{eq:window6_pushforward_markov_spectral_gap} 中的 \(\lambda_\star(P_6)\) 代入计算得到。
\end{proof}

\begin{conclusion}[强可并失败的结构性后果：稳定类型过程在一般初态下不可能为一阶马尔可夫闭合]\label{con:xi-window6-strong-lumpability-failure-structural}
由纤维划分 \(\{F_6^{-1}(x)\}_{x\in X_6}\) 诱导的粗粒化不满足 strong lumpability：
存在同一纤维内两点 \(\omega,\omega'\) 使其一步邻域在稳定类型上的计数分布不同（显式反例见 \eqref{eq:foldbin6_strong_lumpability_counterexample}）。
因此不存在任何仅依赖 \(X_t\) 的齐次转移核 \(P\) 能在\emph{所有}初始微态分布下使 \((X_t)\) 闭合为 \(X_6\) 上的无记忆马尔可夫链。
\end{conclusion}
\begin{proof}[证明要点]
这是命题 \ref{prop:terminal-foldbin6-strong-lumpability-fails}；反例由生成片段 \eqref{eq:foldbin6_strong_lumpability_counterexample} 给出。
\end{proof}

\begin{conclusion}[最小无记忆闭包的强迫性：纤维内瞬时热化与推送核闭合严格等价]\label{con:xi-window6-instant-thermalization-equivalence}
在强可并性失败的背景下，若引入最小的无记忆闭包机制：每一步立方体翻转后，条件于当前稳定类型 \(X_t=x\) 将微态立即均匀化到纤维 \(F_6^{-1}(x)\)（纤维内瞬时热化），
则稳定类型过程闭合为 \(X_6\) 上的可逆马尔可夫链，其转移核必为推送核 \(P_6\)，平衡测度必为 \(\pi_6\propto d^{\mathrm{bin}}_6\)。
等价地，在该协议口径下，一阶无记忆闭合的存在性与“纤维内瞬时热化”严格等价。
\end{conclusion}
\begin{proof}[证明要点]
闭合的唯一性与可逆性由定理 \ref{thm:terminal-foldbin6-pushforward-markov} 与推论 \ref{cor:terminal-foldbin6-minimal-markov-closure-gap} 给出；其强迫性由命题 \ref{prop:terminal-foldbin6-strong-lumpability-fails} 的否定性结论推出。
\end{proof}

\paragraph{Witt--Fourier 双桥与算子化 Big-Witt：统一幽灵标尺与 \texorpdfstring{$p$}{p}-典型护栏}

\begin{conclusion}[Witt--Fourier 双桥的严格粘接：同一角色轴 \texorpdfstring{$\mathrm{Tw}_\chi$}{Tw\_chi} 同时对角化“卷积--碰撞”与“迹--Euler”]\label{con:xi-witt-fourier-commuting-cube}
在有限环 \(G_m\) 及其角色群 \(\widehat G_m\) 上，对每个角色 \(\chi\) 定义纤维侧扭曲
\[
\mathrm{Tw}_\chi(d_m):=\sum_{r\in G_m}d_m(r)\chi(r)=\widehat d_m(\chi),
\]
并在轨道侧以同一 \(\chi\) 定义 Dirichlet 扭曲核 \(T_\chi\)。
则 Fourier 对角化与 Witt/Euler 交换图在同一条 \(\mathrm{Tw}_\chi\) 轴上严格相容，二者沿角色方向粘接形成严格交换立方体；
并且在 \(\Fold_m\) 口径下有完全因子化
\[
\mathrm{Tw}_\chi(d_m)=\prod_{j=1}^m\bigl(1+\chi(F_{j+1})\bigr).
\]
\end{conclusion}
\begin{proof}[证明要点]
交换立方体与扭曲轴一致性见注 \ref{rem:twist-commuting-cube}；
乘积分解见定理 \ref{thm:fold-spectrum-factorization}（并在同一注中给出对齐方式）。
\end{proof}

\begin{conclusion}[统一幽灵标尺：二阶碰撞缺口同时度量 Fourier 非平凡频谱能量与 Witt/Euler 的偏离强度]\label{con:xi-collision-gap-ghost-yardstick}
Witt 侧的“迹 \(\to\) primitive \(\to\) Euler”桥包与 Fourier 侧的 Parseval 能量桥包共享同一个可计算的标量指纹——二阶碰撞缺口（等价于非平凡频谱能量的 \(\ell^2\)-范数）：
\[
\boxed{\ F_{m+2}\,\mathsf{Col}_m-1=\frac{1}{4^m}\sum_{t\ne 0}\bigl|\widehat c_m(t)\bigr|^2\ }.
\]
因此 prime-like（Euler 乘积）侧与谱/迹侧的偏离强度可被压缩到同一个可审计标量上，形成跨通道一致性核验的最短核验回路。
\end{conclusion}
\begin{proof}[证明要点]
见注 \ref{rem:fw-bridge}（并由定理 \ref{thm:fold-collision-spectrum} 的 Parseval 恒等式给出该标尺的严格推导；其与 Rényi-2/均匀性缺口的等价关系见推论 \ref{cor:pom-renyi2-near-uniform}）。
\end{proof}

\begin{conclusion}[碰撞概率的谱迹公式与 Riesz--Parseval 恒等式]\label{con:xi-fold-collision-spectral-trace-parseval}
令 \(M:=F_{m+2}\)，并把折叠多重度剖面视为模环上的函数
\(
d_m:\ZZ/(M\ZZ)\to\NN
\)。
定义归一化分布
\[
\mu_m(r):=\frac{d_m(r)}{2^m},
\qquad
u(r):=\frac{1}{M}\qquad\bigl(r\in\ZZ/(M\ZZ)\bigr),
\]
以及碰撞概率
\[
\mathsf{Col}_m:=\sum_{r\in\ZZ/(M\ZZ)}\mu_m(r)^2.
\]
则有严格谱迹恒等式
\[
\boxed{\;
\mathsf{Col}_m
=\frac{1}{M}\sum_{k\in\ZZ/(M\ZZ)}
\prod_{j=1}^{m}\cos^2\!\Bigl(\pi k\,\frac{F_{j+1}}{M}\Bigr)\; }.
\]
因此
\[
\mathsf{Col}_m-\frac1M
=\frac{1}{M}\sum_{k\ne 0}
\prod_{j=1}^{m}\cos^2\!\Bigl(\pi k\,\frac{F_{j+1}}{M}\Bigr).
\]
并且 \(\chi^2\)-散度与碰撞缺口满足精确等价
\[
\boxed{\;
\chi^2(\mu_m\mid u)
:=\sum_{r}\frac{\bigl(\mu_m(r)-u(r)\bigr)^2}{u(r)}
=M\sum_{r}\Bigl(\mu_m(r)-\frac1M\Bigr)^2
=M\mathsf{Col}_m-1
=M\Bigl(\mathsf{Col}_m-\frac1M\Bigr).\;}
\]
\end{conclusion}
\begin{proof}[证明要点]
由定理 \ref{thm:fold-collision-spectrum} 的 Parseval 恒等式，
\(
\mathsf{Col}_m=\frac{1}{M}\sum_k|\widehat\mu_m(k)|^2
\)。
再由定理 \ref{thm:fold-spectrum-factorization} 的乘积分解得到
\(
|\widehat\mu_m(k)|^2=\prod_{j=1}^m\cos^2(\pi kF_{j+1}/M)
\)，代回即得谱迹公式。
\(\chi^2\) 恒等式则由
\(
M\sum_r(\mu_m(r)-1/M)^2=M\sum_r\mu_m(r)^2-1
\)
直接化简得到。证毕。
\end{proof}

\begin{conclusion}[群环卷积算子的零点--核维数等价]\label{con:xi-fold-group-ring-convolution-kernel-dimension}
在群环
\[
\CC[\ZZ/(M\ZZ)]\cong\CC[x]/(x^M-1)
\]
中令生成函数
\[
D_m(x):=\sum_{r=0}^{M-1}d_m(r)x^r,
\]
并定义卷积/乘法算子
\[
T_m:\CC[x]/(x^M-1)\to\CC[x]/(x^M-1),\qquad f\longmapsto D_m\cdot f\ \ (\mathrm{mod}\ x^M-1).
\]
记谱零点集合
\[
\mathcal Z_m:=\{k\in\ZZ/(M\ZZ):\ \widehat d_m(k)=0\}.
\]
则核维数与零点数严格一致：
\[
\boxed{\;
\dim_\CC\ker(T_m)=|\mathcal Z_m|,
\qquad
\mathrm{rank}(T_m)=M-|\mathcal Z_m|.\;}
\]
\end{conclusion}
\begin{proof}[证明要点]
取标准 Fourier 基（或等价地：把 \(\CC[x]/(x^M-1)\) 以 \(x\mapsto e^{-2\pi i k/M}\) 的特征分解对角化）。
在该基下，乘法算子 \(T_m\) 对角化，其特征值为
\(
D_m(e^{-2\pi i k/M})=\widehat d_m(k)
\)（见命题 \ref{prop:fold-multiplicity-group-algebra} 与定理 \ref{thm:fold-spectrum-factorization}）。
故核由所有 \(\widehat d_m(k)=0\) 的特征方向张成，维数等于零点数目。证毕。
\end{proof}

\begin{conclusion}[谱零点的同余方程刻画与余类分层]\label{con:xi-fold-spectrum-zero-congruence-coset-stratification}
令 \(M:=F_{m+2}\) 且 \(k\in\ZZ/(M\ZZ)\)。
则
\[
k\in\mathcal Z_m
\quad\Longleftrightarrow\quad
\exists j\in\{1,\dots,m\}\ \text{s.t.}\ e^{-2\pi i kF_{j+1}/M}=-1.
\]
等价地（因此必有 \(M\) 偶）：
\[
k\in\mathcal Z_m
\quad\Longleftrightarrow\quad
\bigl(M\ \text{为偶数}\bigr)\ \wedge\
\exists j\in\{1,\dots,m\}\ \text{s.t.}\ kF_{j+1}\equiv \frac{M}{2}\pmod M.
\]
并且对固定的 \(j\)，令 \(g_j:=\gcd(F_{j+1},M)\)。
则同余
\(
kF_{j+1}\equiv M/2\ (\mathrm{mod}\ M)
\)
可解当且仅当 \(g_j\mid M/2\)；
在可解时解集为 \(\ZZ/(M\ZZ)\) 的一个仿射陪集，且其元素个数恰为 \(g_j\)。
\end{conclusion}
\begin{proof}[证明要点]
第一段由定理 \ref{thm:fold-spectrum-factorization}：\(\widehat d_m(k)=\prod_{j=1}^m(1+e^{-2\pi i kF_{j+1}/M})\)，故为零当且仅当某因子为零。
第二段与第三段是线性同余方程的标准判别，并在本文记号下由定理 \ref{thm:fold-zero-coset-rigidity} 给出显式陪集与其大小。证毕。
\end{proof}

\begin{conclusion}[六倍窗的核维数下界与线性不可逆的半阶崩塌]\label{con:xi-fold-window6-kernel-half-index-lower-bound}
若 \(m+2\equiv 0\pmod 6\)，则存在一个半阶余类块满足
\[
|\mathcal Z_m|\ \ge\ F_{(m+2)/2}.
\]
因此由结论 \ref{con:xi-fold-group-ring-convolution-kernel-dimension} 立刻推出
\[
\dim_\CC\ker(T_m)\ge F_{(m+2)/2},
\qquad
\mathrm{rank}(T_m)\le M-F_{(m+2)/2}.
\]
\end{conclusion}
\begin{proof}[证明要点]
谱零点下界为推论 \ref{cor:fold-zero-half-index-multiple6}；
核维数等价由结论 \ref{con:xi-fold-group-ring-convolution-kernel-dimension} 得到。证毕。
\end{proof}

\begin{conclusion}[\texorpdfstring{$\chi^2$}{chi^2} 距离的常数级刚性下界]\label{con:xi-fold-chi2-constant-residual-energy}
令 \(M:=F_{m+2}\)，并沿用结论 \ref{con:xi-fold-collision-spectral-trace-parseval} 的 \(\mu_m,u\)。
设临界共振常数 \(C_\varphi>0\) 如定理 \ref{thm:fold-critical-resonance-constant}。
则有常数级下界
\[
\boxed{\;
\liminf_{m\to\infty}\chi^2(\mu_m\mid u)\ \ge\ 2C_\varphi^2\;}
\]
等价地，
\[
\mathsf{Col}_m-\frac1M\ \ge\ \frac{2C_\varphi^2+o(1)}{M}.
\]
\end{conclusion}
\begin{proof}[证明要点]
由结论 \ref{con:xi-fold-collision-spectral-trace-parseval}，
\(
\chi^2(\mu_m\mid u)=M(\mathsf{Col}_m-1/M)
\)。
而推论 \ref{cor:fold-critical-resonance-gap} 给出
\(
\liminf_{m\to\infty}M(\mathsf{Col}_m-1/M)\ge 2C_\varphi^2
\)。
合并即得。证毕。
\end{proof}

\begin{conclusion}[高阶碰撞核 Perron 根的端点斜率由最大纤维指数锁定]\label{con:xi-fold-high-q-perron-endpoint-slope-maxfiber}
记最大纤维 \(M_m:=\max_{x\in X_m}d_m(x)\)，并令 \(\alpha_m:=\frac1m\log M_m\)。
若极限 \(\alpha_\ast:=\lim_{m\to\infty}\alpha_m\) 存在且有限，则对每个固定整数 \(q\ge 1\)，令
\[
S_q(m):=\sum_{x\in X_m}d_m(x)^q,
\qquad
\log r_q:=\lim_{m\to\infty}\frac{1}{m}\log S_q(m)
\quad(\text{若极限存在}),
\]
则恒有双边夹逼
\[
q\alpha_\ast\ \le\ \log r_q\ \le\ q\alpha_\ast+\log\varphi,
\]
并且端点斜率满足
\[
\lim_{q\to\infty}\frac{1}{q}\log r_q=\alpha_\ast.
\]
\end{conclusion}
\begin{proof}[证明要点]
由平凡夹逼 \(M_m^q\le S_q(m)\le |X_m|\,M_m^q\) 得
\[
q\alpha_m\ \le\ \frac{1}{m}\log S_q(m)\ \le\ q\alpha_m+\frac{1}{m}\log|X_m|.
\]
利用 \(|X_m|=F_{m+2}=\Theta(\varphi^m)\) 与 \(\alpha_m\to\alpha_\ast\) 取极限得到 \(\log r_q\) 的夹逼；
再令 \(q\to\infty\) 得端点斜率。证毕。
\end{proof}

\begin{conclusion}[Rényi 指纹率的零温极限]\label{con:xi-fold-renyi-fingerprint-rate-zero-temp}
沿用结论 \ref{con:xi-fold-high-q-perron-endpoint-slope-maxfiber} 的 \(r_q\) 与 \(\alpha_\ast\)。
定义每阶 \(q\) 的指纹率
\[
h_R(q):=\frac{q\log 2-\log r_q}{q-1}.
\]
则存在零温极限
\[
\boxed{\;\lim_{q\to\infty}h_R(q)=\log 2-\alpha_\ast.\;}
\]
\end{conclusion}
\begin{proof}[证明要点]
由结论 \ref{con:xi-fold-high-q-perron-endpoint-slope-maxfiber}，有 \(\log r_q=q\alpha_\ast+O(1)\)，代入定义并令 \(q\to\infty\) 得结论。证毕。
\end{proof}

\begin{conclusion}[规范体积密度的余维一常数：离散体积与折叠缺陷熵的稳定平移]\label{con:xi-fold-gauge-volume-codim1-constant}
对任意折叠系统，令 \(d_m(w)\) 为纤维大小函数并定义规范群体积
\[
|G_m|:=\prod_{w\in X_m} d_m(w)!,
\qquad
g_m:=\frac{1}{2^m}\log|G_m|,
\qquad
\kappa_m:=\frac{1}{2^m}\sum_{w\in X_m} d_m(w)\log d_m(w).
\]
在计数守恒 \(\sum_w d_m(w)=2^m\) 下，Stirling 展开给出普适差常数：存在与 \(m\) 无关的绝对常数使
\[
g_m=\kappa_m-1+O\!\Bigl(\frac{|X_m|\,m}{2^m}\Bigr).
\]
在本文给定的 Fibonacci 稳定类型尺度 \(|X_m|=F_{m+2}=\Theta(\varphi^m)\) 下，误差为指数级小量，从而
\[
\kappa_m-g_m=1+O\!\bigl(m(\varphi/2)^m\bigr)=1+o(1).
\]
\end{conclusion}
\begin{proof}[证明要点]
见定理 \ref{thm:fold-bin-gauge-volume-stirling-second-order}（并结合命题 \ref{prop:fold-bin-gauge-volume} 的记号设定与余项界）。证毕。
\end{proof}
\paragraph{规范体积密度与输出熵的定量镜像}
在 bin-fold 口径下，纤维信息项 \(\kappa_m\)、输出 Shannon 熵与规范体积密度 \(g_m\) 并非三条彼此独立的复杂度账本，而是由同一条带有显式常数项的 Stirling--entropy 恒等式锁定。

\begin{theorem}[规范体积密度与输出 Shannon 熵之间的一纳特定律]\label{thm:xi-foldbin-gauge-entropy-one-nat-law}
\leanverified{paper\_xi\_foldbin\_gauge\_entropy\_one\_nat\_law}
令
\[
\mu_m(w):=\frac{d_m^{\mathrm{bin}}(w)}{2^m},
\qquad
w\in X_m,
\]
并记
\[
\kappa_m:=\mathbb E_{\mu_m}\!\bigl[\log d_m^{\mathrm{bin}}(W_m)\bigr],
\qquad
g_m:=2^{-m}\log|\mathcal G_m^{\mathrm{bin}}|,
\]
其中 \(W_m\sim \mu_m\)。
则有精确恒等式
\[
\kappa_m=m\log 2-H(\mu_m),
\]
其中
\[
H(\mu_m):=-\sum_{w\in X_m}\mu_m(w)\log\mu_m(w).
\]
进一步，
\[
g_m
=
m\log 2-H(\mu_m)-1
+\frac{(\varphi/2)^m}{2\sqrt5}
\Bigl(
\varphi^2 m\log\frac{2}{\varphi}
+\varphi^2\log(2\pi)-\log\varphi
\Bigr)
+O\!\bigl((\varphi/2)^m\bigr),
\]
从而
\[
m\log 2-H(\mu_m)-g_m\longrightarrow 1.
\]
\end{theorem}
\begin{proof}
由 \(\mu_m(w)=d_m^{\mathrm{bin}}(w)/2^m\) 得
\[
\log d_m^{\mathrm{bin}}(w)=m\log 2+\log\mu_m(w).
\]
对 \(\mu_m\) 取期望可得
\[
\kappa_m
=
\sum_{w\in X_m}\mu_m(w)\log d_m^{\mathrm{bin}}(w)
=
m\log 2+\sum_{w\in X_m}\mu_m(w)\log\mu_m(w)
=
m\log 2-H(\mu_m).
\]

另一方面，命题 \ref{prop:fold-bin-gauge-density-two-scale} 给出
\[
g_m=\kappa_m-1
+\frac{(\varphi/2)^m}{2\sqrt5}
\Bigl(
\varphi^2 m\log\frac{2}{\varphi}
+\varphi^2\log(2\pi)-\log\varphi
\Bigr)
+O\!\bigl((\varphi/2)^m\bigr).
\]
将 \(\kappa_m=m\log 2-H(\mu_m)\) 代入即得所示展开。
最后，由结论 \ref{con:xi-fold-gauge-volume-codim1-constant} 已知
\(
\kappa_m-g_m=1+O(m(\varphi/2)^m)
\)，故
\[
m\log 2-H(\mu_m)-g_m
=
\kappa_m-g_m
\longrightarrow 1.
\]
证毕。
\end{proof}

\begin{corollary}[bin-fold 输出 Shannon 熵与 gauge 密度的显式常数项镜像]\label{cor:xi-foldbin-shannon-gauge-constant-mirror}
\leanverified{paper\_xi\_foldbin\_shannon\_gauge\_constant\_mirror}
沿用定理 \ref{thm:xi-foldbin-gauge-entropy-one-nat-law} 的记号。则有展开
\[
\kappa_m
=
m\log\frac{2}{\varphi}
-\frac{\log\varphi}{1+\varphi^2}
+O\!\bigl((\varphi/2)^m\bigr)
=
m\log\frac{2}{\varphi}
-\frac{\log\varphi}{\varphi\sqrt5}
+O\!\bigl((\varphi/2)^m\bigr),
\]
从而
\[
H(\mu_m)
=
m\log\varphi
+\frac{\log\varphi}{\varphi\sqrt5}
+O\!\bigl((\varphi/2)^m\bigr),
\]
以及
\[
g_m
=
m\log\frac{2}{\varphi}
-1
-\frac{\log\varphi}{\varphi\sqrt5}
+O\!\bigl(m(\varphi/2)^m\bigr).
\]
\end{corollary}
\begin{proof}
当 \(q=1\) 时，
\(
\pi_{m,1}(w)=d_m^{\mathrm{bin}}(w)/S_1^{\mathrm{bin}}(m)=d_m^{\mathrm{bin}}(w)/2^m=\mu_m(w)
\)，
故命题 \ref{prop:fold-bin-escort-lastbit} 在 \(q=1\) 时给出
\[
\kappa_m
=
m\log\frac{2}{\varphi}
-\frac{\log\varphi}{1+\varphi^2}
+O\!\bigl((\varphi/2)^m\bigr).
\]
由恒等式 \(1+\varphi^2=\varphi\sqrt5\) 得第二种写法。

再由定理 \ref{thm:xi-foldbin-gauge-entropy-one-nat-law}，
\[
H(\mu_m)=m\log 2-\kappa_m.
\]
代入上式即得 Shannon 熵展开。

最后，由结论 \ref{con:xi-fold-gauge-volume-codim1-constant}，
\[
g_m=\kappa_m-1+O\!\bigl(m(\varphi/2)^m\bigr),
\]
将 \(\kappa_m\) 的展开代入即得 \(g_m\) 的常数项公式。证毕。
\end{proof}


\begin{conclusion}[Big-Witt 运算的算子编译：原子 Euler 数据的半环可同态嵌入迹类算子组合]\label{con:xi-atomic-witt-operator-compilation}
令 \(\mathcal{W}_{\mathrm{atom}}\) 为由原子数据 \((\beta,n,m)\) 生成的 Big-Witt 型半环，令 \(\mathrm{TC}_1\) 为迹类算子在酉等价类下的交换半环（加法为直和、乘法为张量积）。
则映射
\[
(\beta,n,m)\longmapsto \bigoplus_{i=1}^{m} C(n,\alpha),\qquad \alpha^n=\beta,
\]
延拓为构造性的半环同态
\[
\boxed{\ \mathcal{W}_{\mathrm{atom}}\ \longrightarrow\ \mathrm{TC}_1\ }.
\]
因此 Big-Witt 的抽象 Euler/plethystic 运算可被编译为迹类算子的直和/张量积组合运算，并与 Fredholm \(\zeta\) 接口严格相容，形成可执行的 Euler 代数。
\end{conclusion}
\begin{proof}[证明要点]
见定理 \ref{thm:atomic-witt-into-tc1}（乘法的 \(\gcd/\mathrm{lcm}\) 规则由命题 \ref{prop:cyclic-block-tensor-gcd-lcm} 保证）。
\end{proof}

\begin{conclusion}[\texorpdfstring{$p$}{p}-典型 Frobenius--Verschiebung 的迹层实现：\texorpdfstring{$\mathsf F_p\circ\mathsf V_p$}{Fp o Vp} 等价于“直和乘以 \texorpdfstring{$p$}{p}”]\label{con:xi-ptypical-frobenius-verschiebung-trace}
在 \(p\)-典型分量上，Frobenius 与 Verschiebung 的算子实现满足 ghost 兼容性与 Verschiebung--Frobenius 关系：
对任意由 \(p\)-幂循环块有限直和得到的算子 \(L\)，对一切 \(n\ge 1\) 有
\[
\boxed{\ \Tr\!\bigl(\mathsf F_p(L)^n\bigr)=\Tr\!\bigl(L^{pn}\bigr)\ },
\qquad
\boxed{\ \mathsf F_p\!\bigl(\mathsf V_p(L)\bigr)\ \simeq\ L^{\oplus p}\ }.
\]
该关系把 Dwork 型同余从“多项式/序列”提升为“迹类算子组合”的结构恒等式，使 \(p\)-进护栏以算子函子方式零成本继承。
\end{conclusion}
\begin{proof}[证明要点]
见命题 \ref{prop:cyclic-p-typical-frobenius-verschiebung}。
\end{proof}

\paragraph{模障碍、cyclotomic 因子与 Smith 形：共振角色表的有限终止审计}

\begin{conclusion}[模障碍与 cyclotomic 因子的完全等价：共振来源归约为单位根因子分类]\label{con:xi-mod-obstruction-cyclotomic-equivalence}
对 Ihara/动力学 \(\zeta\) 的扭曲通道 \(\det(I-tB_\rho(u))^{-1}\)，存在非平凡模障碍当且仅当某个 \(\rho\neq \mathbf{1}\) 使相应谱半径达到平凡通道；
且该条件当且仅当 \(\det(I-tB_\rho(u))\) 含 cyclotomic 因子（归一化后在单位圆上出现单位根因子）。
因此“共振来源”可被等价归约为有限维整系数多项式的 cyclotomic 因子分类问题。
\end{conclusion}
\begin{proof}[证明要点]
见结论 \ref{con:xi-mod-obstruction-cyclotomic-factor} 及其所引分类定理 \ref{thm:conclusion75-ihara-cyclotomic-factor-classification}。
\end{proof}

\begin{conclusion}[障碍群的同调分类与多项式时间可计算性：全共振角色表可编译为一次 SNF 计算]\label{con:xi-resonant-character-snf-computable}
共振角色集合 \(X_{\mathrm{res}}\) 构成有限阿贝尔群，且可由同调商的 torsion 部分刻画；
其不变因子分解与生成角色可由相应整数矩阵的 Smith 正规形直接读出。
更强地，存在确定性算法在时间 \(\mathrm{poly}(|V|,|E|,\log W)\) 内输出障碍不变量与 \(X_{\mathrm{res}}\) 的不变因子分解及生成角色，从而“是否存在 cyclotomic 因子/共振通道”的全角色审计可在有限步内终止。
\end{conclusion}
\begin{proof}[证明要点]
同调刻画见定理 \ref{thm:gm-resonant-character-group-homology}；Smith 形与多项式时间可计算性见命题 \ref{prop:gm-smith-compute-obstructions}。
\end{proof}

\subsubsection{中心六重聚合与 Perron 分支的六分支惯性：解析半径与代数次数约束}

\begin{conclusion}[中心六重聚合的几何后果：中心唯一实交与最近复奇点距离给出解析半径]\label{con:xi-rate-center-sixfold-unique-real-rstar}
同步核速率曲线的规范化消元多项式 \(R(\alpha,u)\) 具有对偶自对称结构（注 \ref{rem:rate-algebraic-elimination-structure}），并在中心截面满足分解
\[
R\!\left(\tfrac12,u\right)=-\frac{(u-1)^6}{64}\,Q(u).
\]
其中 \(Q(u)\) 系数回文，故 \(Q(u)=u^{16}Q(1/u)\)，零点以互逆对出现（并与复共轭合并为四元组）；并且 Sturm 计数给出 \(Q\) 在实轴无零点，从而 \(Q(u)>0\) 对所有 \(u\in\RR\) 成立（命题 \ref{prop:rate-center-Q-reciprocal-no-real}）。
因此 \(R(\tfrac12,u)=0\) 在实参数上当且仅当 \(u=1\)，且重数恰为 \(6\)（推论 \ref{cor:rate-center-unique-real}）。
进一步定义最近复奇点距离
\[
R_\star:=\min\{|u-1|:\ Q(u)=0\},
\]
则 \(R_\star\) 给出中心分支的解析半径并可数值审计（推论 \ref{cor:rate-center-Rstar}）。
\end{conclusion}
\begin{proof}[证明要点]
自对称与中心分解的可复算证书见注 \ref{rem:rate-algebraic-elimination-structure}；
\(Q\) 的自倒易性与实根空性见命题 \ref{prop:rate-center-Q-reciprocal-no-real}；
中心唯一实交与六重性见推论 \ref{cor:rate-center-unique-real}；
最近复奇点距离的数值证书见推论 \ref{cor:rate-center-Rstar}。
\end{proof}

\begin{conclusion}[惯性群强迫的“六分支单元”：Perron 分支代数次数必被 \texorpdfstring{$6$}{6} 整除]\label{con:xi-rate-center-perron-degree-multiple-6}
沿用中心六重聚合证书 \(R(\tfrac12,u)=-(u-1)^6Q(u)/64\)，并假设一般横截条件 \(\partial_\alpha R(\tfrac12,1)\neq 0\)。
令 \(F(\alpha,u)\in\QQ[\alpha,u]\) 为 \(R(\alpha,u)\) 在 \(\QQ(\alpha)[u]\) 中包含点 \((\tfrac12,1)\) 的不可约因子（例如包含 Perron 分支的那一因子）。
则在函数域扩张 \(\QQ(\alpha)\subset \QQ(\alpha,u)\ (F(\alpha,u)=0)\) 中，\(\alpha=\tfrac12\) 处的分歧指数为 \(e=6\)，从而 \(\deg_uF\)（亦即该分支对 \(\QQ(\alpha)\) 的代数次数）必被 \(6\) 整除。
等价地，中心点邻域的 Perron 分支具有不可约的六分支惯性结构，任何试图以单值解析支“穿越中心点”的方案在代数层面被排除。
\end{conclusion}
\begin{proof}[证明要点]
见推论 \ref{cor:rate-center-perron-degree-multiple-6}；惯性群语义对齐见注 \ref{rem:rate-center-inertia-c6}。
\end{proof}

\paragraph{中心六重临界分歧向有限节点反演病态的传递}
中心六重聚合把参数扰动压缩到六次根尺度之后，Prony/Hankel 型有限窗口反演的不适定性会沿同一尺度被放大为显式幂律爆炸。

\begin{theorem}[六重临界分歧强制 \texorpdfstring{$\kappa$}{kappa}-节点谱反演的普适爆炸率]\label{thm:xi-center-sixfold-kappa-node-inversion-blowup}
沿用结论 \ref{con:xi-rate-center-perron-degree-multiple-6} 的记号，并在 \((\alpha,u)=(\tfrac12,1)\) 的邻域内取包含 Perron 分支的局部解析支。
假设存在常数 \(c_1,c_2>0\) 使得当 \(\alpha\to \tfrac12\) 时有
\[
c_1\left|\alpha-\frac12\right|^{1/6}
\le
|u-1|
\le
c_2\left|\alpha-\frac12\right|^{1/6}.
\]

固定 \(\kappa\ge 2\)，并考虑由长度 \(2\kappa-1\) 的 Prony/Hankel 数据恢复 \(\kappa\) 个谱节点
\[
r_1(u),\dots,r_\kappa(u)
\]
的局部反演问题。
记最小分离度为
\[
\operatorname{sep}(u):=\min_{j\neq k}|r_j(u)-r_k(u)|.
\]
若各成对节点在 \(u=1\) 处都横截裂分，即对每个 \(j\neq k\) 都存在 \(c_{jk}\neq 0\) 使
\[
r_j(u)-r_k(u)=c_{jk}(u-1)+O\bigl((u-1)^2\bigr),
\]
则任意局部 Lipschitz 右逆的 Lipschitz 常数 \(L(\alpha)\) 都满足下界
\[
L(\alpha)\ge C_\kappa\left|\alpha-\frac12\right|^{-(\kappa-1)/6},
\]
其中 \(C_\kappa>0\) 仅依赖于 \(\kappa\)、采样规范以及局部横截常数。
\end{theorem}
\begin{proof}
由各成对差值的横截裂分，存在常数 \(c_3,c_4>0\) 使当 \(u\) 充分接近 \(1\) 时，
\[
c_3|u-1|
\le
\operatorname{sep}(u)
\le
c_4|u-1|.
\]
再与假设的六重临界分歧尺度比较合并，得到
\[
\operatorname{sep}(u)\asymp \left|\alpha-\frac12\right|^{1/6}.
\]

另一方面，命题 \ref{con:xi-collision-illposedness-sep-blowup} 已给出任意基于长度 \(2\kappa-1\) 有限样本的 Prony/Hankel 型局部反演都满足
\[
L\ge C'_\kappa\,\operatorname{sep}(u)^{-(\kappa-1)},
\]
其中 \(C'_\kappa>0\) 仅依赖于 \(\kappa\) 与采样规范。
将上式中的 \(\operatorname{sep}(u)\) 以 \(|\alpha-\tfrac12|^{1/6}\) 代换，并把隐含比较常数吸收入新的常数 \(C_\kappa\)，即得
\[
L(\alpha)\ge C_\kappa\left|\alpha-\frac12\right|^{-(\kappa-1)/6}.
\]
证毕。
\end{proof}
\endinput


\endinput
